\documentclass{zhvt-classic}

% 使用本地源流明体（GenRyuMin2 TW）按文件路径加载，不依赖系统字体名
% 版框（内框／界栏／书口侧栏／鱼尾）改用与版心印章同一印泥棕色：
%   印章图片 yshs.png 实色 = RGB(135,68,52) = HTML 874434
\usepackage{xcolor}
\usepackage{tikz}
% xcolor 与 l3color 是两套色彩注册表，需各注册同名 sealbrown：
%   xcolor 版 → 供 TikZ/pgf 的 fill= 读取；L3 版 → 供 cls 版框 \color_stroke:n 读取。
\definecolor{sealbrown}{HTML}{874434}
\ExplSyntaxOn
\color_set:nnn { sealbrown } { HTML } { 874434 }
\ExplSyntaxOff

\zhvtset{
  grid color       = sealbrown,
  main font        = GenRyuMin2TW-SB,
  main font path   = C:/Users/Administrator/texmf/fonts/opentype/public/GenRyuMin2TW/,
jiazhu font      = GenRyuMin2TW-M,
  jiazhu font path = C:/Users/Administrator/texmf/fonts/opentype/public/GenRyuMin2TW/,
  % 页码框容纳 4 字：正常页「九十七」+前置空 1 格=4 行恰好；
  % 100+ 页「一百零一」共 5 行超出，末字「一/二/三」被框裁掉（2026-10-04 实测），故放宽到 5
  shukou page number chars = 5,
  % 书口书名框高从默认 6 字减到 4 字，把省出的 2 字高拨给 chapter 框，
  % 使 6 字书名+卷名放进 9 字高的 chapter 框后不再压线（2026-10-05）
  shukou title chars = 4,
}

\title{史记}[夏本纪]
\author{汉·司马迁}
\maker{闲叟漫记·重製}

% ===== 索隐述赞样式（2026-11-08 愚大人定稿参数，与《五帝本纪》完全一致）=====
%   题头「索隐述赞曰」五字：黑体 + 圆角棕框反白。「曰」是《索隐》述赞的引导语，
%   并入题头后，正文首组直接起四言，五组字数齐整（每组 16 格）。
%   正文韵文：仿宋 \fontsize{24}{41} —— 列距 41pt 对齐全书界行网格 121+41k，字不压线。
%   ⚠ 自定义字体家族必须 RawFeature={vertical:+vert} 仅 +vert；
%     simhei 若加 +vhal 会整体加载失败、静默回退成仿宋。
\newCJKfontfamily{\shuzanfont}[Path=C:/Windows/Fonts/,RawFeature={vertical:+vert}]{simhei.ttf}
\newCJKfontfamily{\shuzanbody}[Path=C:/Windows/Fonts/,RawFeature={vertical:+vert}]{simfang.ttf}
%   定尺寸框 152×30pt（五字 23pt 横排约 115pt，左右各留约 18pt），圆角 6pt；
%   baseline={(0,2)} 使整框落入界行格子、不压内细框线；
%   白字节点 (0pt,8.2pt)：x 不补偿（simhei 竖排墨迹水平本就居中），y 补 8.2pt。
%   ⚠ 只画一层白字——叠「隐形」棕字节点会在框外露淡影，视觉上就是「字超出底纹」。
\newcommand{\shuzanhead}[1]{%
\begin{tikzpicture}[baseline={(0,2)}]
\fill[sealbrown, rounded corners=6pt]
(-76pt,-15pt) rectangle (76pt,15pt);
\node[text=white, inner sep=0pt, anchor=center] at (0pt,8.2pt)
{\shuzanfont\fontsize{23}{23}\selectfont #1};
\end{tikzpicture}%
}

% ★ 修正书口卷名位置：chapter 框起始位置改回 cls 默认值，避免全局影响。
%   只在每个卷首页通过 \juan 命令里的空命令微调该页位置（安全，不会跨页污染）。
\ExplSyntaxOn
\AtBeginDocument{
  \dim_set:Nn \l__zhvt_shukou_chapter_title_start_dim
    { \l__zhvt_shukou_title_height_dim + \l__zhvt_font_size_dim * 2 }
  \dim_set:Nn \l__zhvt_shukou_page_number_start_dim
    { \l__zhvt_shukou_chapter_title_start_dim + \l__zhvt_shukou_chapter_title_height_dim }
}
\ExplSyntaxOff

\begin{document}

% ★★ 版心（书口）排布覆写：把书名放到鱼尾之下（2026-10-04）
%
%   cls 默认版心自上而下：书名（title 静态框，在鱼尾之上）→ 鱼尾 → 篇名（chapter 动态框）→ 页码
%   本模板改为：          鱼尾（位置不动）→ 书名 → 篇名 → 页码分割线 → 页码
%
%   手法：把 cls 放在 title 静态框里的书名，改放进 chapter 动态框（鱼尾之下）；
%         title 静态框留空。书名取 \l__zhvt_title_str，即 \title{...} 的第一参数，
%         所以换书时只改 \title，不必在此重复写书名。
%
%   ⚠ 三条必须遵守的写法（2026-10-04 实测）：
%     ① 汉字必须用 \hbox{...} 包住。flowfram 会把框内容装进一个宽 14.4mm（= 一个 line_sep
%        = 书口栏宽）的纵向 minipage 并按该宽度折行，而整页靠 PDF /Rotate 90 显示、
%        框内文字本身不变向 —— 不包 \hbox 时「围炉夜话」4 字会被逐字折成 4 行，
%        在版面上就成了横排 4 字并横向溢出版框。
%     ② 竖排下**一个 \hbox 就是一整列**。书名与篇名要同处一列，必须放进同一个 \hbox；
%        分成两个 \hbox 时，后一个会被当成新的一列、甩到隔壁栏去
%        （实测「序」跑到书口栏外侧 14.4mm 的正文列上）。中间的 \quad 在框内是横向空白，
%        显示出来正好是书名与篇名之间的纵向间隔。
%     ③ ★ 内容盒要写成 **与框等宽的定宽盒** `\hbox to \l__zhvt_shukou_chapter_title_height_dim{…\hfil}`。
%        实测：chapter 框是把内容盒沿纵向**居中**摆放的，内容宽度一变，书名就跟着漂：
%        自然宽度的 \hbox（书名 38.4mm）落在 x=157.2mm，加篇名后（57.6mm）落在 x=152.4mm，
%        目录页（无篇名）与正文页的书名就差半格。做成等宽盒后，无论有没有篇名，
%        内容固定从框顶起排，实测两页均落在 x=142.8mm，不再漂。
\ExplSyntaxOn
% 把书口 chapter 框（书名+篇名）起始位置下移 1 字高，避免 4 字卷名顶到鱼尾
% ⚠ cls 的 \__zhvt_calculate_dimensions: 在 \AtBeginDocument 里执行，会重算该 dim；
%   必须把覆写挂到 begin document 之后，否则会被二次覆盖（2026-10-05 实测）。
% (旧方案已移到导言区，直接修改 chapter 框起始位置)
% 卷标题命令 #1=书口篇名(鱼尾) #2=正文卷名（含卷序，卷名内空格用 ~）
%
% ★★ 目录页两项修订（2026-10-05，愚大人指出）★★
%   ① 页码起点统一对齐：
%      原状 —— 目录项写成 \par\hspace{4em}卷名\hspace{4em}页码，页码起点 = 卷名实宽 + 4em。
%      卷名实宽随字数变：「卷一~梁惠王上」6 字 +~、「卷七~离娄上」5 字 +~，
%      差 1 个字位 → 卷七~卷十 4 条的页码整体比其余 10 条高 27.21pt（一个字），看着上下不齐。
%      改法 —— 卷名补全角空格（宽 = 1 字 = 27.215pt）补足到 **7 字符位**再写进目录，
%      14 条页码起点遂全部落在同一条竖线上（起点对齐；末字随页码字数自然参差）。
%      ⚠ 只对目录里的副本 \l_tmpc_tl 补，正文卷名仍用 \l_tmpa_tl 原样 —— 正文像素不受影响。
%   ② 卷名整体上移半字位：前导 \hspace 由 4em 改 3.5em。
%      实测原状目录块在版框内偏下（天头留白 139.6pt / 地脚留白 110.7pt），
%      上移 0.5em = 13.6pt 后上下基本均分，视觉居中。
\NewDocumentCommand { \juan } { mm } {
	\clearpage
	\tl_set:Nn \l__zhvt_chapter_title_tl { #1 }
	\tl_set:Nn \l_tmpa_tl { #2 }
	\tl_set_eq:NN \l_tmpc_tl \l_tmpa_tl
	\int_step_inline:nn { 7 - \str_count:n { #2 } }
	  { \tl_put_right:Nn \l_tmpc_tl { 　 } }
	\tl_set:Nx \l_tmpb_tl { \zhnum { page } }
	\tl_put_left:Nn \l_tmpb_tl { \hspace { 4em } }
	\tl_put_left:Nn \l_tmpb_tl { \l_tmpc_tl }
	\tl_put_left:Nn \l_tmpb_tl { \par \hspace { 3.5em } }
	\addtocontents { toc } { \l_tmpb_tl }
	\skip_horizontal:n { \l__zhvt_font_size_dim * 4 }
	\l_tmpa_tl
	\par
	\thispagestyle { empty }
}
\cs_set_protected:Npn \__zhvt_set_shukou_content: {
	% ★ 鱼尾上方留白（title 静态框）嵌入闲叟钤印：
	%   整页 /Rotate 90，故此处把图用 \rotatebox{90} 转正为竖章；
	%   尺寸卡在版心条幅内（书内约 7mm 宽 × 33.8mm 高，留余量不压鱼尾）。
	\setstaticcontents * { title } {
		\vbox_to_ht:nn { \l__zhvt_baseline_skip_dim } {
			\vskip 14.2pt
			\hbox_to_wd:nn { \l__zhvt_shukou_title_height_dim } {
				\hfil
				\rotatebox [ origin = c ] { 90 } { \includegraphics [ width = 7mm , height = 33.8mm ] { yshs.png } }
				\hfil
			}
		}
	}
% ★ 版心（书口）统一改法（2026-11-16 愚大人定）：18pt＋框内居中＋书口文字同版框色
%   ① 原 27.2pt 时 7 字长篇名（如「廉颇蔺相如列传」）末字溢出 chapter 框、压住页码；
%      改 \fontsize{18pt}{18pt}，前后 \hfil 令字串在框内居中（不随篇名字数漂移）。
%   ② \vbox{\vskip 8.2pt} 定量下移，抵消缩字号后栏宽方向 8.2pt 上偏（像素实测 y 中心 543.3 = 原版）。
%      ⚠ \raisebox / \vspace* 在 flowfram 框内无效；换字号须同步换补偿：16pt→10.1、19pt→7.4。
%   ③ chapter 与 page~number 均置 \color{sealbrown}（HTML 874434），与版框／界行同色；正文仍黑。
%   ⚠ 页码字位、钤印、鱼尾、界行、正文一律未动；五部合订本同此口径。
	\setdynamiccontents * { chapter }
	{ \color{sealbrown} \vbox { \vskip 8.2pt \hbox to \l__zhvt_shukou_chapter_title_height_dim
		{ \hfil \fontsize{18pt}{18pt}\selectfont \l__zhvt_title_str \quad \l__zhvt_chapter_title_tl \hfil } } }
	\setdynamiccontents * { page~number } { \color{sealbrown} \quad \zhnum { page } }
}
\__zhvt_set_shukou_content:
\ExplSyntaxOff

\maketitle

% 单篇，无多卷目录，不渲染目录页




% 方案①：仅改 A 类错讹，头部与版式不变

\juan{夏本纪}{卷二~夏本纪第二}
\hspace*{2em}夏禹。名曰文命。禹之父曰鲧。鲧之父曰帝颛顼。颛顼之父曰昌意。昌意之父曰黄帝。禹者。黄帝之玄孙而帝颛顼之孙也。禹之曾大父昌意及父鲧皆不得在帝位。为人臣。\\
\hspace*{2em}当帝尧之时。鸿水滔天。浩浩怀山襄陵。下民其忧。尧求能治水者。群臣四岳皆曰鲧可。尧曰。﹁鲧为人负命毁族。不可。﹂四岳曰。﹁等之未有贤于鲧者。愿帝试之。﹂于是尧听四岳。用鲧治水。九年而水不息。功用不成。于是帝尧乃求人。更得舜。舜登用。摄行天子之政。巡狩。行视鲧之治水无状。乃殛鲧于羽山以死。天下皆以舜之诛为是。于是舜举鲧子禹。而使续鲧之业。\\
\hspace*{2em}尧崩。帝舜问四岳曰。﹁有能成美尧之事者使居官。﹂皆曰。﹁伯禹为司空。可成美尧之功。﹂舜曰。﹁嗟。然。﹂命禹。﹁女平水土。维是勉之。﹂禹拜稽首。让于契。后稷。皋陶。舜曰。﹁女其往视尔事矣。﹂\\
\hspace*{2em}禹为人敏给克勤。其德不违。其仁可亲。其言可信。声为律。身为度。称以出。亹亹穆穆。为纲为纪。\\
\hspace*{2em}禹乃遂与益。后稷奉帝命。命诸侯百姓兴人徒以傅土。行山表木。定高山大川。禹伤先人父鲧功之不成受诛。乃劳身焦思。居外十三年。过家门不敢入。薄衣食。致孝于鬼神。卑宫室。致费于沟淢。陆行乘车。水行乘船。泥行乘橇。山行乘檋。左准绳。右规矩。载四时。以开九州。通九道。陂九泽。度九山。令益予众庶稻。可种卑湿。命后稷予众庶难得之食。食少。调有余相给。以均诸侯。禹乃行相地宜所有以贡。及山川之便利。\\
\hspace*{2em}禹行自冀州始。冀州。既载壶口。治梁及岐。既修太原。至于岳阳。覃怀致功。致于衡漳。其土白壤。赋上上错。田中中。常。卫既从。大陆既为。鸟夷皮服。夹右碣石。入于海。\\
\hspace*{2em}济。河维沇州。九河既道。雷夏既泽。雍。沮会同。桑土既蚕。于是民得下丘居土。其土黑坟。草繇木条。田中下。赋贞。作十有三年乃同。其贡漆。丝。其篚织文。浮于济。漯。通于河。\\
\hspace*{2em}海岱维青州。堣夷既略。潍。淄其道。其土白坟。海滨广潟。厥田斥卤。田上下。赋中上。厥贡盐絺。海物维错。岱畎丝。枲。铅。松。怪石。莱夷为牧。其篚酓丝。浮于汶。通于济。\\
\hspace*{2em}海岱及淮维徐州。淮。沂其治。蒙。羽其艺。大野既都。东原厎平。其土赤埴坟。草木渐包。其田上中。赋中中。贡维土五色。羽畎夏狄。峄阳孤桐。泗滨浮磬。淮夷蠙珠臮鱼。其篚玄纤缟。浮于淮。泗。通于河。\\
\hspace*{2em}淮海维扬州。彭蠡既都。阳鸟所居。三江既入。震泽致定。竹箭既布。其草惟夭。其木惟乔。其土涂泥。田下下。赋下上上杂。贡金三品。瑶。琨。竹箭。齿。革。羽。旄。岛夷卉服。其篚织贝。其包橘。柚锡贡。均江海。通淮。泗。\\
\hspace*{2em}荆及衡阳维荆州。江。汉朝宗于海。九江甚中。沱。涔已道。云土。梦为治。其土涂泥。田下中。赋上下。贡羽。旄。齿。革。金三品。杶。榦。栝。柏。砺。砥。砮。丹。维箘簬。楛。三国致贡其名。包匦菁茅。其篚玄纁玑组。九江入赐大龟。浮于江。沱。涔。于汉。逾于雒。至于南河。\\
\hspace*{2em}荆。河惟豫州。伊。雒。瀍。涧既入于河。荥播既都。道荷泽。被明都。其土壤。下土坟垆。田中上。赋杂上中。贡漆。丝。絺。纻。其篚纤絮。锡贡磬错。浮于雒。达于河。\\
\hspace*{2em}华阳黑水惟梁州。汶。嶓既艺。沱。涔既道。蔡。蒙旅平。和夷厎绩。其土青骊。田下上。赋下中三错。贡璆。铁。银。镂。砮。磬。熊。罴。狐。貍。织皮。西倾因桓是来。浮于潜。逾于沔。入于渭。乱于河。\\
\hspace*{2em}黑水西河惟雍州。弱水既西。泾属渭汭漆。沮既从。沣水所同。荆。岐已旅。终南。敦物至于鸟鼠。原隰厎绩。至于都野。三危既度。三苗大序。其土黄壤。田上上。赋中下。贡璆。琳。琅玕浮于积石。至于龙门西河。会于渭汭织皮昆仑。析支。渠搜。西戎即序。\\
\hspace*{2em}道九山。汧及岐至荆山。逾于河。壶口。雷首至于太岳。砥柱。析城至于王屋。太行。常山至于碣石。入于海。西倾。朱圉。鸟鼠至于太华。熊耳。外方。桐柏至于负尾。道嶓冢。至于荆山。内方至于大别。汶山之阳至衡山。过九江。至于敷浅原。\\
\hspace*{2em}道九川。弱水至于合黎。余波入于流沙。道黑水。至于三危。入于南海。道河积石。至于龙门。南至华阴。东至砥柱。又东至于盟津。东过雒汭。至于大邳。北过降水。至于大陆。北播为九河。同为逆河。入于海。嶓冢道漾。东流为汉。又东为苍浪之水。过三澨。入于大别。南入于江。东汇泽为彭蠡。东为北江。入于海。汶山道江。东别为沱。又东至于醴。过九江。至于东陵。东迆北会于汇。东为中江。入于海。道沇水。东为济。入于河。泆为荥。东出陶丘北。又东至于荷。又东北会于汶。又东北入于海。道淮自桐柏。东会于泗。沂。东入于海。道渭自鸟鼠同穴。东会于沣。又东北至于泾。东过漆。沮。入于河。道雒自熊耳。东北会于涧。瀍。又东会于伊。东北入于河。\\
\hspace*{2em}于是九州攸同。四奥既居。九山刊旅。九川涤原。九泽既陂。四海会同。六府甚脩。众土交正。致慎财赋。咸则三壤成赋。中国赐土姓。﹁祗台德先。不距朕行。﹂\\
\hspace*{2em}令天子之国以外五百里甸服。百里赋纳总。二百里纳铚。三百里纳秸服。四百里粟。五百里米。甸服外五百里侯服。百里采。二百里任国。三百里诸侯。侯服外五百里绥服。三百里揆文教。二百里奋武卫。绥服外五百里要服。三百里夷。二百里蔡。要服外五百里荒服。三百里蛮。二百里流。\\
\hspace*{2em}东渐于海。西被于流沙。朔。南暨。声教讫于四海。于是帝锡禹玄圭。以告成功于天下。天下于是太平治。\\
\hspace*{2em}皋陶作士以理民。帝舜朝。禹。伯夷。皋陶相与语帝前。皋陶述其谋曰。﹁信其道德。谋明辅和。﹂禹曰。﹁然。如何。﹂皋陶曰。﹁於。慎其身修。思长。敦序九族。众明高翼。近可远在已。﹂禹拜美言。曰。﹁然。﹂皋陶曰。﹁於。在知人。在安民。﹂禹曰。﹁吁。皆若是。惟帝其难之。知人则智。能官人。能安民则惠。黎民怀之。能知能惠。何忧乎驩兜。何迁乎有苗。何畏乎巧言善色佞人。﹂皋陶曰。﹁然。於。亦行有九德。亦言其有德。﹂乃言曰。﹁始事事。宽而栗。柔而立。愿而共。治而敬。扰而毅。直而温。简而廉。刚而实。强而义。章其有常吉哉。日宣三德。蚤夜翊明有家。日严振敬六德。亮采有国。翕受普施。九德咸事。俊乂在官。百吏肃谨。毋教邪淫奇谋。非其人居其官。是谓乱天事。天讨有罪。五刑五用哉。吾言厎可行乎。﹂禹曰。﹁女言致可绩行。﹂皋陶曰。﹁余未有知。思赞道哉。﹂\\
\hspace*{2em}帝舜谓禹曰。﹁女亦昌言。﹂禹拜曰。﹁於。予何言。予思日孳孳。﹂皋陶难禹曰。﹁何谓孳孳。﹂禹曰。﹁鸿水滔天。浩浩怀山襄陵。下民皆服于水。予陆行乘车。水行乘舟。泥行乘橇。山行乘檋。行山刊木。与益予众庶稻鲜食。以决九川致四海。浚畎浍致之川。与稷予众庶难得之食。食少。调有余补不足。徙居。众民乃定。万国为治。﹂皋陶曰。﹁然。此而美也。﹂\\
\hspace*{2em}禹曰。﹁於。帝。慎乃在位。安尔止。辅德。天下大应。清意以昭待上帝命。天其重命用休。﹂帝曰。﹁吁。臣哉。臣哉。臣作朕股肱耳目。予欲左右有民。女辅之。余欲观古人之象。日月星辰。作文绣服色。女明之。予欲闻六律五声八音。来始滑。以出入五言。女听。予即辟。女匡拂予。女无面谀。退而谤予。敬四辅臣。诸众谗嬖臣。君德诚施皆清矣。﹂禹曰。﹁然。帝即不时。布同善恶则毋功。﹂\\
\hspace*{2em}帝曰。﹁毋若丹朱傲。维慢游是好。毋水行舟。朋淫于家。用绝其世。予不能顺是。﹂禹曰。﹁予娶涂山。辛壬癸甲。生启予不子。以故能成水土功。辅成五服。至于五千里。州十二师。外薄四海。咸建五长。各道有功。苗顽不即功。帝其念哉。﹂帝曰。﹁道吾德。乃女功序之也。﹂\\
\hspace*{2em}皋陶于是敬禹之德。令民皆则禹。不如言。刑从之。舜德大明。\\
\hspace*{2em}于是夔行乐。祖考至。群后相让。鸟兽翔舞。箫韶九成。凤凰来仪。百兽率舞。百官信谐。帝用此作歌曰。﹁陟天之命。维时维几。﹂乃歌曰。﹁股肱喜哉。元首起哉。百工熙哉。﹂皋陶拜手稽首扬言曰。﹁念哉。率为兴事。慎乃宪。敬哉。﹂乃更为歌曰。﹁元首明哉。股肱良哉。庶事康哉。﹂舜又歌曰。﹁元首丛脞哉。股肱惰哉。万事堕哉。﹂帝拜曰。﹁然。往钦哉。﹂于是天下皆宗禹之明度数声乐。为山川神主。\\
\hspace*{2em}帝舜荐禹于天。为嗣。十七年而帝舜崩。三年丧毕。禹辞辟舜之子商均于阳城。天下诸侯皆去商均而朝禹。禹于是遂即天子位。南面朝天下。国号曰夏后。姓姒氏。\\
\hspace*{2em}帝禹立而举皋陶荐之。且授政焉。而皋陶卒。封皋陶之后于英。六。或在许。而后举益。任之政。\\
\hspace*{2em}十年。帝禹东巡狩。至于会稽而崩。以天下授益。三年之丧毕。益让帝禹之子启。而辟居箕山之阳。禹子启贤。天下属意焉。及禹崩。虽授益。益之佐禹日浅。天下未洽。故诸侯皆去益而朝启。曰﹁吾君帝禹之子也。﹂于是启遂即天子之位。是为夏后帝启。\\
\hspace*{2em}夏后帝启。禹之子。其母涂山氏之女也。\\
\hspace*{2em}有扈氏不服。启伐之。大战于甘。将战。作甘誓。乃召六卿申之。启曰。﹁嗟。六事之人。予誓告女。有扈氏威侮五行。怠弃三正。天用剿绝其命。今予维共行天之罚。左不攻于左。右不攻于右。女不共命。御非其马之政。女不共命。用命。赏于祖。不用命。僇于社。予则帑僇女。﹂遂灭有扈氏。天下咸朝。\\
\hspace*{2em}夏后帝启崩。子帝太康立。帝太康失国。昆弟五人。须于洛汭。作五子之歌。\\
\hspace*{2em}太康崩。弟中康立。是为帝中康。帝中康时。羲。和湎淫。废时乱日。胤往征之。作胤征。\\
\hspace*{2em}中康崩。子帝相立。帝相崩。子帝少康立。帝少康崩。子帝予立。帝予崩。子帝槐立。帝槐崩。子帝芒立。帝芒崩。子帝泄立。帝泄崩。子帝不降立。帝不降崩。弟帝扃立。帝扃崩。子帝廑立。帝廑崩。立帝不降之子孔甲。是为帝孔甲。帝孔甲立。好方鬼神。事淫乱。夏后氏德衰。诸侯畔之。天降龙二。有雌雄。孔甲不能食。未得豢龙氏。陶唐既衰。其后有刘累。学扰龙于豢龙氏。以事孔甲。孔甲赐之姓曰御龙氏。受豕韦之后。龙一雌死。以食夏后。夏后使求。惧而迁去。\\
\hspace*{2em}孔甲崩。子帝皋立。帝皋崩。子帝发立。帝发崩。子帝履癸立。是为桀。帝桀之时。自孔甲以来而诸侯多畔夏。桀不务德而武伤百姓。百姓弗堪。乃召汤而囚之夏台。已而释之。汤修德。诸侯皆归汤。汤遂率兵伐夏桀。桀走鸣条。遂放而死。桀谓人曰。﹁吾悔不遂杀汤于夏台。使至此。﹂汤乃践天子位。代夏朝天下。汤封夏之后。至周封于杞也。\\
\hspace*{2em}太史公曰。禹为姒姓。其后分封。用国为姓。故有夏后氏。有扈氏。有男氏。斟寻氏。彤城氏。褒氏。费氏。杞氏。缯氏。辛氏。冥氏。斟氏。戈氏。孔子正夏时。学者多传夏小正云。自虞。夏时。贡赋备矣。或言禹会诸侯江南。计功而崩。因葬焉。命曰会稽。会稽者。会计也。\\
% ===== 索隐述赞（司马贞《史记索隐》卷二十九·十二本纪述赞·夏本纪）=====
%   每四句一列：\hbox to \hsize{\hfil 组文 \hfil} 定宽盒内 \hfil 居中、字距不拉伸；
%   盒间 \linebreak 供断点（★ \hspace*{30em} 强制换列会吞掉下一列首字，绝不可用）。
%   \clearpage 让题头 + 5 组共 6 列同页完整呈现，且主文列距保持 41pt 不被压缩。
\clearpage
\shuzanhead{索隐述赞曰}\linebreak
{\shuzanbody\fontsize{24}{41}\selectfont
\hbox to \hsize{\hfil 尧遭鸿水。黎人阻饥。禹勤沟洫。手足胼胝。\hfil}\linebreak
\hbox to \hsize{\hfil 言乘四载。动履四时。娶妻有日。过门不私。\hfil}\linebreak
\hbox to \hsize{\hfil 九土既理。玄圭锡兹。帝启嗣立。有扈违命。\hfil}\linebreak
\hbox to \hsize{\hfil 五子作歌。太康失政。羿浞斯侮。夏室不竞。\hfil}\linebreak
\hbox to \hsize{\hfil 降于孔甲。扰龙乖性。嗟彼鸣条。其终不令。\hfil}%
\par}
\end{document}
